% -- Encoding UTF-8 without BOM
% -- XeLaTeX => PDF (BIBER)

\documentclass[german, print]{cv-style-wide}

\usepackage{changepage}

\usepackage[unicode,
            colorlinks=false,
            bookmarks=true,
            bookmarksopen=true,
            bookmarksdepth=2
           ]{hyperref}
\hypersetup{
    pdftitle={Curriculum Vitae},
    pdfauthor={},
    pdfsubject={CV},
    colorlinks=false
}

\begin{document}

\settowidth{\bwidth}{BSc}
\settowidth{\mwidth}{MSc}
\settowidth{\pwidth}{PhD}

\header{Vorname}{Nachname}{}

\begin{center}
home: \href{https://example.com}{example.com} \quad
github: anonymized \quad
linkedin: anonymized \quad
xing: anonymized \quad
\href{mailto:email@example.com}{email@example.com} \\
Stadt, Land \quad Staatsangehörigkeit, Wohnsitz \\
\end{center}

\lastupdated

\vspace{0.5cm}

Innovativer angewandter Mathematiker mit langjähriger Erfahrung in wissenschaftlicher Datenverarbeitung und über 15 Jahren \\
Berufserfahrung in der angewandten Forschung mit Schwerpunkt auf medizinischen Interventionen. International anerkannter Experte \\
für Ultraschallsimulation in Therapie und Diagnostik. Auf der Suche nach neuen Herausforderungen in der Modellierung zur \\
Unterstützung der Entwicklung medizinischer Geräte. Umzugsbereit.
\begin{NoHyphItemize}
  \item Erstellt realistische und gleichzeitig handhabbare mathematische Modelle sowie einsetzbare, getestete und dokumentierte Simulationen. Ausgezeichnete Programmier- und Softwareentwicklungskenntnisse, erfüllt ISO 13485/IEC 62304.
  \item Kommunikationsfähigkeiten, verfeinert durch Arbeit in multidisziplinären, internationalen Teams, Lehr- und Vortragstätigkeit an Universitäten und Konferenzen.
  \item 17 begutachtete Zeitschriftenartikel (über 250 Zitierungen), ein Buchkapitel, eingeworbene Fördermittel (über 500.000 €), 12 eingeladene Vorträge, Betreuung von drei Doktoranden und einem Masterstudenten sowie Pflege des Open-Source-Codes \href{https://k-wave-python.readthedocs.io}{k-wave-python}.
\end{NoHyphItemize}

\section{Berufliche Erfahrungen}

\begin{entrylist}

\entry
  {Nov 2019--}
  {Fraunhofer-Institut für Digitale Medizin MEVIS}
  {Bremen, Deutschland}
  {\jobtitle{Leitender Wissenschaftler — Modellierungs- und Simulationsgruppe\\
   \textnormal{Fähigkeiten: python, VTK, Ultraschallmodellierung, Behandlungsplanung, Softwareentwicklung, Elastographie, Unsicherheitsanalyse}}%
   \begin{itemize}[topsep=0pt]
     \item Entwicklung groß angelegter Simulationen für Mikrowellen- und Ultraschallablationstherapien durch Konzeption und Parallelisierung leistungsstarker numerischer Methoden zur klinisch relevanten Behandlungsplanung.
     \item Entwicklung schneller Ultraschall-Beamforming-Algorithmen und eines transkraniellen akustischen/elastischen Ausbreitungssimulators mittels GPU-beschleunigter Techniken.
   \end{itemize}}

\entry
  {Jun 2014--Nov 2019}
  {National Physical Laboratory}
  {Teddington, Vereinigtes Königreich}
  {\jobtitle{Leitender Wissenschaftler — Gruppe für Ultraschall und Unterwasserakustik\\
   \textnormal{Fähigkeiten: python, matlab, FEM (COMSOL, FeniCS), Ultraschallmodellierung, Signalverarbeitung}}%
   \begin{itemize}
     \item Etablierung messbasierter Simulation nichtlinearer Ausbreitung komplexer Medien; Integration empirischer Daten in Rechenmodelle, Aufnahme in IEC 63587.
   \end{itemize}}

\entry
  {Jun 2011--Jun 2014}
  {Institute of Cancer Research / The Royal Marsden Hospital}
  {Sutton, Vereinigtes Königreich}
  {\jobtitle{Postdoktoranden-Wissenschaftler — Therapeutische Ultraschallgruppe / Physik\\
   \textnormal{Fähigkeiten: python, VTK, ultrasound modelling, treatment planning, software development}}%
   \begin{itemize}
     \item Entwicklung eines hochintensiven fokussierten Ultraschall-Behandlungsplanungssystems mit Multi-Element-Phased-Array zur präzisen Fokussteuerung und Sicherheit.
   \end{itemize}}

\entry
  {Jun 2008--Jun 2011}
  {University College London}
  {London, Vereinigtes Königreich}
  {\jobtitle{Postdoktoranden-Wissenschaftler — Ultraschallgruppe / Maschinenbau\\
   \textnormal{Fähigkeiten: Mathematische Modellierung, Fortran, Differentialgleichungen}}%
   \begin{itemize}
     \item Untersuchung des Einflusses von Kavitation auf therapeutischen Ultraschall mittels numerischer und analytischer Verfahren zur Optimierung von Wirksamkeit und Sicherheit.
   \end{itemize}}

\end{entrylist}

\section{Ausbildung}

\begin{entrylist}
\threeentry{2004--2008}{PhD — Dynamical Systems}{University College London, Vereinigtes Königreich}{}
\threeentry{2003--2004}{MSc — Modern Applications of Mathematics}{University of Bath, Vereinigtes Königreich}{}
\threeentry{2000--2003}{BSc — Mathematics with Applied Math./Math. Physics}{Imperial College London, Vereinigtes Königreich}{}
\end{entrylist}

\section{Auszeichnungen \& Wertschätzungsindikatoren}

\begin{entrylist}
\entry
  {2020}
  {IEEE IUS Challenge on Ultrasound Beamforming with Deep Learning (CUBDL)}
  {}
  {Gemeinsamer erster Platz bei internationalem ML-Wettbewerb zur Ultraschallbildrekonstruktion (2020)}

\entry
  {2015--}
  {Internationaler Experte}
  {}
  {Mitglied IEC/BSI TC 87 (Ultrasonics), Mitarbeiter Institute of Mathematics and its Applications, Vollmitglied Institute of Physics}

\entry
  {Mehrere}
  {Erweiterte Stipendien}
  {}
  {Fördermittel für MSc (2004), PhD (2008) und Postdoc (2014) durch EPSRC.}
\end{entrylist}

\section{Zusatzqualifikationen}

\begin{tabular}{p{0.6\textwidth} p{0.4\textwidth}}
\textbf{Programming:} python, C++, Matlab, OpenCL & \textbf{Libraries:} ITK, VTK, boost, eigen \\
\textbf{DevOps:} git, svn, github, gitlab, google test, pytest, make, cmake, visual studio & \textbf{Computation:} FEniCS, Comsol \\
\textbf{Sprachen:} Englisch (Muttersprache), Deutsch (B2.1) & \\
\end{tabular}

\end{document}